% Ad Atticum 7.7 (ad-atticum-07-07) --- English translation
% Translated by Claude (Anthropic) from the Latin (Perseus).
% Style guide: STYLE.md.

\ciceroLetterOpener{CICERO ATTICO salutem}

\ciceroSection{1} Dionysius, an excellent man --- as I too have found him by now --- and most learned and most devoted to you, came to Rome on the fifteenth day before the Kalends of January and delivered me your letter. For there are so many words about Dionysius in your letter, but \textit{[corrupt]} you do not add that he gave thanks. And yet certainly he ought to have done so; and if he had, then --- such is your humanity --- you would have added it. To me, however, no \textit{[Greek: recantation]} regarding him is allowed, on account of the testimony of my earlier letter. Let him be, then, by all means a good man. He has done well by this very fact, that he has given me the further opportunity of getting to know him through and through. Philogenes wrote correctly to you: he saw to what he owed. I wished him to make use of that money as long as he might; and so he has used it for forty-three months.

\ciceroSection{3} For Pomptinus I want recovery, and as to what you write, that he has entered the City, I fear what it may mean; for he would not have done that without grave cause. As for myself, since the 4th day before the Nones of January is the day of the \textit{Compitalia}, I do not wish to come to the Alban villa on that day, so as not to come as a burden to the household. So, on the third day before the Nones of January; thence to the City on the day before the Nones. On what day your \textit{[Greek: attack]} of fever falls I do not know, but I do not want you, on any account, to stir for the inconvenience of your health.

\ciceroSection{4} About my honor --- unless Caesar shall have privately stirred up something through his tribunes --- the rest seems quiet enough; and most quiet of all is my own mind, which takes the whole thing in good part, and all the more because I now hear from many quarters that it has been settled by Pompey and his council to send me into Sicily, on the strength of the imperium I hold. This is \textit{[Greek: Abderite logic]}. For neither has the Senate decreed nor the People ordered that I should hold imperium in Sicily. But if the commonwealth refers this matter to Pompey, why send me rather than some private citizen? Accordingly, if this imperium is to be a burden to me, I shall use the first exit I see. As for what you write, that there is a marvelous expectation of me, and that yet no one of the loyalists, or of those tolerably so, doubts what I shall do --- I do not understand whom you call ``loyalists.'' I myself know of none, that is, if it is whole orders of loyalists we are looking for: as individuals there are good men. But in our dissensions one has to look for orders and classes of the loyal. Do you call the Senate loyal, by whose doing the provinces are without governors with imperium (for Curio could never have held out, if measures had begun to be taken against him; the Senate refused to follow that line; from which it has come about that no successor is being sent to Caesar)? --- or the publicans, who have never been firm and now are most friendly to Caesar? --- or the financiers? --- or the farmers, whose dearest wish is quiet? --- unless you suppose those men afraid of falling under a despotism, who, so long as they were left in their quiet, never refused it. What, then?

\ciceroSection{6} ``Do you approve of his command being prolonged after the day of the law has passed, while he is keeping his army?'' I do not even approve of it in absence; but when that has once been granted, the other was granted along with it. For if a ten-year command, granted in the very form it was, is approved, then it follows that my own banishment is approved, and the loss of the Campanian land, and the adoption of a patrician by a plebeian, of a man of Gades by a man of Mytilene, and the wealth of Labienus, and the wealth of Mamurra, and the gardens of Balbus, and his Tusculan estate. But of all these there is one source. The weak should have been resisted, and that was easy; now there are eleven legions, as much cavalry as he shall wish, the Transpadanes, the urban plebs, so many tribunes of the plebs, a youth so utterly corrupted, a commander of such authority, of such audacity. With this man one must either fight it out, or let his command be taken into account under the law.

\ciceroSection{7} ``Fight it out,'' you say, ``rather than be a slave.'' To what end? If you are beaten, you will be proscribed; if you win, you will still be a slave. ``What, then,'' you say, ``will you do?'' What cattle do, which when driven asunder follow the herds of their kind. As the ox follows the herd, so I shall follow the loyal men, or whoever shall be called loyal, even if they crash. What is best, given how badly matters have been brought to this pass, I see plainly. For no one knows for certain, once it comes to arms, what will be; but this all know, that if the loyalists are defeated, this man will be no more merciful in the slaughter of leading men than Cinna was, nor more moderate in the property of the rich than Sulla was. \textit{[Greek: I have been doing politics with you]} for some time now, and would do it longer, did not my lamp fail me. To sum up: ``Say your piece, Marcus Tullius.'' I assent to Gnaeus Pompeius, that is, to Titus Pomponius. Please give my greetings to Alexis, that most charming of boys --- unless perhaps in my absence he has grown into a young man, which he seemed to be in the way of doing.
